%--------------------------------------------------------------------%
% BAB V PENGUJIAN DAN EVALUASI
%--------------------------------------------------------------------%
\chapter{PENGUJIAN DAN EVALUASI} \label{bab:5}
\petunjuk{Hasil pengujian dibandingkan dengan rencana pada 2.4 dan keterpenuhan tujuan proyek. Utamakan pengujian bagian Anda dan kontribusinya pada keterpenuhan tujuan proyek. Bagian ini wajib ditulis sendiri dan berbeda dari laporan anggota tim lain (Pedoman TA 5.1.3). \tahap{\tidakperlu{Tidak diperlukan}}{\wajib{WAJIB}}}

\section{Hasil Uji} \label{sec:5.1}
\petunjuk{Hasil tiap uji dibandingkan dengan kriteria lulus pada rencana, beserta statusnya. Data hasil uji yang banyak boleh disajikan dalam tabel bernomor.}
\tulisuraian

\section{Verifikasi TKT 3} \label{sec:5.2}
\petunjuk{Bukti bahwa TKT minimal 3 tercapai dan siapa yang memvalidasi (Pedoman TA butir 3.3.2).}
\tulisuraian

\section{Demonstrasi Produk} \label{sec:5.3}
\petunjuk{Skenario demonstrasi langsung yang sesuai dengan klaim pada laporan.}
\tulisuraian

\section{Keterpenuhan Tujuan dan Kriteria Penerimaan} \label{sec:5.4}
\petunjuk{Cocokkan tiap tujuan dan kriteria penerimaan pada 1.3 dengan hasil dan nyatakan statusnya.}
\tulisuraian
